\begin{lemma}
\label{lemma-functorial-triangles-decomposable}
Let $\mathcal{D}$ be a pre-triangulated category. Let
$\operatorname{Ar}(\mathcal{D})$ be the category whose objects are morphisms
of $\mathcal{D}$ and whose morphisms are commutative squares. Let
$\operatorname{DTri}(\mathcal{D})$ be the category whose objects are
distinguished triangles and whose morphisms are morphisms of triangles. Let
$$
U : \operatorname{DTri}(\mathcal{D}) \longrightarrow
\operatorname{Ar}(\mathcal{D})
$$
be the functor which sends a distinguished triangle to its first arrow.
Assume $U$ has a section. If $\mathcal{D}$ has countable products or countable
coproducts, then every morphism of $\mathcal{D}$ is isomorphic to a composition
$$
K \oplus Z \longrightarrow Z \longrightarrow Z \oplus Q
$$
of a projection and a coprojection. In particular, $\mathcal{D}$ has kernels
and cokernels.
\end{lemma}

\begin{proof}
By Lemma \ref{lemma-projectors-have-images-triangulated}, the category
$\mathcal{D}$ is Karoubian. Let $T$ be a section of $U$, let $f : X \to Y$,
and write
$$
T(f) = (X, Y, Z_f, f, g_f, h_f).
$$
Set $d_f = -h_f[-1] : Z_f[-1] \to X$. Exactness of the sequence obtained from
this triangle gives a natural epimorphism
$$
q_f : \Hom_\mathcal{D}(-, Z_f[-1]) \longrightarrow \mathcal{K}_f,
\qquad
\mathcal{K}_f(W) = \{v : W \to X \mid f \circ v = 0\},
$$
given by composition with $d_f$.

Write $T(W \to 0) = (W, 0, Z_W, 0, 0, h_W)$. By two rotations and
Lemma \ref{lemma-third-object-zero}, the map $h_W : Z_W \to W[1]$ is an
isomorphism. Given $v \in \mathcal{K}_f(W)$, applying $T$ to the commutative
square $(v, 0) : (W \to 0) \to (X \to Y)$ gives a morphism of triangles
$(v, 0, c_v)$. Define
$$
s_{f,W}(v) = -c_v[-1] \circ h_W[-1]^{-1}.
$$
Functoriality of $T$ shows that these maps are natural in $W$. The last square
in the morphism of triangles gives
$h_f[-1] \circ c_v[-1] = v \circ h_W[-1]$; hence the chosen signs give
$q_f \circ s_f = 1_{\mathcal{K}_f}$.

Thus $e_f = s_f \circ q_f$ is an idempotent natural endomorphism of
$\Hom_\mathcal{D}(-, Z_f[-1])$. By Yoneda it comes from an idempotent
$e : Z_f[-1] \to Z_f[-1]$. Since $\mathcal{D}$ is Karoubian, $e$ splits, and
the image of this splitting represents $\mathcal{K}_f$. Thus $f$ has a kernel.
The conclusion now follows from Lemma \ref{lemma-when-split}.
\end{proof}

